The extra item-draw experiment separates a fixed design's behavior from changes in source membership or clustering. At $k=100$, residual stratification has lower exact design variance than uniform difference on \RandomLargeVarianceWins/13 random splits and \UpperLargeVarianceWins/13 upper-tail splits, but only \LowerLargeVarianceWins/13 lower-tail splits (Figure~\ref{fig:variance}). The ratios of macro-averaged exact variances are \RandomLargeVarianceRatio, \UpperLargeVarianceRatio, and \LowerLargeVarianceRatio, respectively. At $k=20$ those ratios rise to \RandomSmallVarianceRatio, \UpperSmallVarianceRatio, and \LowerSmallVarianceRatio: the fixed ten-stratum design is inefficient at the smallest budget. These exact variances are diagnostic calculations on full held-out vectors, never inputs to sampling.

Nominal normal $95$\% intervals for residual strata cover \RandomResidualStrataLargeCoverage\%, \UpperResidualStrataLargeCoverage\%, and \LowerResidualStrataLargeCoverage\% at $k=100$ in the three protocols. At $k=20$ coverage falls to \RandomResidualStrataSmallCoverage\%, \UpperResidualStrataSmallCoverage\%, and \LowerResidualStrataSmallCoverage\%. Uniform difference has better small-budget coverage (\RandomDifferenceSmallCoverage\%, \UpperDifferenceSmallCoverage\%, and \LowerDifferenceSmallCoverage\%), although it too falls below nominal. The appendix reports interval widths and benchmark-block intervals. Design-unbiased point and variance estimation therefore should not be described as an exact finite-sample confidence procedure.
